\section{Conclusion}
LEAF treats zero-shot multi-target defense as a problem of constructing reusable threat context. A shared attention operator refines current entities, a current-self query reads the intermediate stages, and learned fusion supplies the resulting feature to downstream temporal memory. This separates within-decision scene processing from inference across physical time.

Across the evaluated SkyHAD populations, LEAF achieves lower mean target damage than the learned baselines and Rule. Historical design comparisons favor current-state reading, while frozen interventions show dependence on constructed context; the effects of sharing, fixed-depth training, and fusion remain configuration-dependent. Frozen execution remains effective over one to six rounds, while normal reads emphasize early context. Higher interception and lower target-damage incidence connect these representations to the intended protection outcome. The structural propositions explain permutation consistency and sufficient local action stability; population transfer and closed-loop depth robustness remain empirical findings within the tested setting.
